\documentclass[11pt]{article}
\usepackage[letterpaper,margin=0.68in]{geometry}
\usepackage{amsmath,amssymb}
\usepackage{enumitem}
\usepackage{xcolor}
\usepackage{fancyhdr}
\usepackage{microtype}
\usepackage{tabularx,array,booktabs}
\usepackage{tikz}
\usepackage[most]{tcolorbox}

\definecolor{olive}{HTML}{4D5430}
\definecolor{gold}{HTML}{9A7B22}
\definecolor{pale}{HTML}{F5F1DC}
\definecolor{softgreen}{HTML}{EDF0E2}
\definecolor{softgray}{HTML}{F5F5F2}
\definecolor{ink}{HTML}{292D20}

\pagestyle{fancy}
\fancyhf{}
\lhead{\textcolor{olive}{MATH\& 107: Math in Society}}
\rhead{\textcolor{gold}{Fall 2026}}
\cfoot{\thepage}
\setlength{\headheight}{14pt}
\setlength{\parindent}{0pt}
\setlength{\parskip}{5pt}
\setlist[enumerate]{leftmargin=*,itemsep=7pt,topsep=5pt}
\setlist[itemize]{leftmargin=1.35em,itemsep=3pt,topsep=3pt}

\newtcolorbox{keybox}[1]{colback=pale,colframe=olive,boxrule=0.8pt,arc=2mm,title=\textbf{#1},fonttitle=\color{olive},left=2mm,right=2mm,top=1.5mm,bottom=1.5mm}
\newtcolorbox{examplebox}[1]{colback=softgreen,colframe=gold,boxrule=0.8pt,arc=2mm,title=\textbf{#1},fonttitle=\color{ink},left=2mm,right=2mm,top=1.5mm,bottom=1.5mm}
\newtcolorbox{refbox}[1]{colback=softgray,colframe=gold,boxrule=0.7pt,arc=2mm,title=\textbf{#1},fonttitle=\color{ink},left=2mm,right=2mm,top=1.5mm,bottom=1.5mm}

\newcommand{\summaryhead}[2]{%
  \clearpage
  {\Large\bfseries\textcolor{olive}{Module #1: #2 — Summary}}\par
  \vspace{2pt}\textcolor{gold}{\rule{\linewidth}{1.2pt}}\par\smallskip
}
\newcommand{\prequizhead}[2]{%
  \clearpage
  {\Large\bfseries\textcolor{olive}{Module #1: #2 — Prequiz}}\par
  \textit{Open-note. Show your mathematical work. A calculation and a clearly stated final answer are enough unless a sentence is specifically requested.}\par
  Name: \hspace{2.6in} Date: \hspace{1.1in}\par
  \vspace{2pt}\textcolor{gold}{\rule{\linewidth}{1.2pt}}\par\smallskip
}
\newcommand{\quizhead}[2]{%
  \clearpage
  {\Large\bfseries\textcolor{olive}{Module #1: #2 — Matching Quiz}}\par
  \textit{These problems use the same skills as the prequiz, with new values and contexts. Show enough work to make your reasoning visible.}\par
  Name: \hspace{2.6in} Date: \hspace{1.1in}\par
  \vspace{2pt}\textcolor{gold}{\rule{\linewidth}{1.2pt}}\par\smallskip
}
\newcommand{\worksmall}{\par\vspace{0.38in}}
\newcommand{\workmed}{\par\vspace{0.62in}}
\newcommand{\worklarge}{\par\vspace{0.92in}}
\newcommand{\workxl}{\par\vspace{1.22in}}
\newcommand{\choice}[1]{\(\square\)~#1\qquad}

\begin{document}
\summaryhead{5}{Basic Computer and Information Theory}

\textbf{What this module is about.} Computers store representations. Binary notation uses powers of \(2\) instead of powers of \(10\), but the place-value idea is the same. Some fractions can be represented exactly in finitely many binary digits and others require approximation. Data can be compressed when it contains predictable repetition, and extra redundant bits can be added when error detection matters. Information theory measures the information in an outcome using its probability.

\begin{keybox}{Place value in base two}
The binary places are
\[
\ldots,2^4,2^3,2^2,2^1,2^0
\]
for whole numbers and
\[
2^{-1},2^{-2},2^{-3},\ldots
\]
for binary fractions.

For example,
\[
1101_2=1(8)+1(4)+0(2)+1(1)=13_{10}.
\]
A byte contains \(8\) bits, so an unsigned byte can represent the integers \(0\) through \(255\).
\end{keybox}

\begin{examplebox}{Example 1: Converting a decimal number to binary}
To represent \(18\), write it as a sum of powers of \(2\):
\[
18=16+2=1(16)+0(8)+0(4)+1(2)+0(1).
\]
Therefore
\[
18_{10}=10010_2.
\]
The symbols changed, but the numerical value did not.
\end{examplebox}

\begin{keybox}{Binary fractions and exactness}
A reduced rational number has a terminating binary expansion exactly when its denominator is a power of \(2\). For example,
\[
\frac{3}{16}=\frac{3}{2^4}
\]
terminates in binary, while \(3/10\) does not because the reduced denominator contains a factor of \(5\).

Example:
\[
0.101_2=1\left(\frac12\right)+0\left(\frac14\right)+1\left(\frac18\right)=\frac58=0.625.
\]
\end{keybox}

\begin{examplebox}{Example 2: Error detection with even parity}
Suppose the data bits are \(1011\). There are three \(1\)s, which is odd. To make the total number of \(1\)s even, append parity bit \(1\):
\[
1011\boxed{1}.
\]
If one bit flips during transmission, the receiver will see an odd number of \(1\)s and can detect that something went wrong. Parity does not correct the message and does not detect every possible multi-bit error; it is a simple error-detection device.
\end{examplebox}

\begin{keybox}{Compression and information}
Run-length encoding replaces a long repeated run by a count and a symbol. For example,
\[
\text{HHHHHTTTTT}\longrightarrow 5\text{H}5\text{T}.
\]
It works well when repetition is common and poorly when symbols change constantly.

The information content of an outcome with probability \(p\) is
\[
I=-\log_2(p).
\]
If \(p=1/32=2^{-5}\), then \(I=5\) bits. Each time probability is halved, information increases by \(1\) bit.
\end{keybox}

\prequizhead{5}{Basic Computer and Information Theory}
\begin{refbox}{Useful facts}
Binary place values are powers of \(2\). A reduced fraction terminates in binary only when its denominator is a power of \(2\). Even parity makes the total number of \(1\)s even. Information: \(I=-\log_2 p\).
\end{refbox}

\begin{enumerate}
\item A sensor reports operating mode \(13\). The controller stores the mode in binary.
\begin{enumerate}[label=(\alph*),itemsep=4pt]
\item Convert \(13_{10}\) to binary.
\item Convert the received binary value \(10110_2\) back to base ten.
\end{enumerate}
\worklarge

\item Convert the binary fraction \(0.101_2\) to a base-ten decimal. Then decide whether each rational number has a finite binary expansion: \(3/16\) and \(3/10\). State the denominator feature you used.
\worklarge

\item A transmitter uses even parity.
\begin{enumerate}[label=(\alph*),itemsep=4pt]
\item Data bits are \(1011\). What parity bit should be appended?
\item A receiver gets the five-bit packet \(10110\). Does the parity check detect an error? Explain by counting the \(1\)s.
\end{enumerate}
\worklarge

\item Apply run-length encoding to the string
\[
\text{HHHHHTTTTT}.
\]
Then explain why the string \(\text{XYZXYZXYZ}\) is not well suited to run-length encoding.
\worklarge

\item An alarm event has probability \(1/32\). How many bits of information does the event carry according to \(I=-\log_2(p)\)? Show the power-of-two step that makes the calculation immediate.
\workmed
\end{enumerate}

\quizhead{5}{Basic Computer and Information Theory}
\begin{enumerate}
\item Convert \(22_{10}\) to binary, and convert \(10001_2\) to base ten.
\workmed

\item Convert \(0.011_2\) to base ten. Decide whether \(5/8\) and \(1/5\) have finite binary expansions.
\workmed

\item A transmitter uses even parity. Data bits are \(0111\). Determine the parity bit. Then decide whether the received packet \(11011\) passes the even-parity check.
\workmed

\item Run-length encode \(\text{MMMMMMMMMM}\). State whether the encoded form is shorter than writing all ten symbols.
\workmed

\item One of \(128\) equally likely outcomes occurs. What is the probability of a particular outcome, and how many bits of information does observing it provide?
\workmed
\end{enumerate}

% =====================================================================
% MODULE 6
% =====================================================================
\end{document}
